\subsection{Training by next-token prediction}
\label{sec:transformer-training}

Training begins with complete token sequences. For
\begin{center}
  \texttt{[BOS, I, like, tea, EOS]},
\end{center}
the input at each position predicts the token immediately to its right. This creates four
targets: \texttt{I}, \texttt{like}, \texttt{tea}, and \texttt{EOS}. The model factorizes
the sequence probability as
\begin{equation}
  p_\theta(x_1,\ldots,x_T)
  =\prod_{t=1}^{T}p_\theta(x_t\mid x_{<t}),
\end{equation}
and minimizes the mean negative log-probability of the observed next tokens.

All correct tokens are already known during training. We therefore feed the whole
sequence at once, while the causal mask prevents position $t$ from seeing a future
target. This procedure is teacher forcing. The resulting logits have shape $[B,T,V]$:
$B$ independent sequences, $T$ next-token decisions per sequence, and $V$ candidate
tokens per decision. The token dimension cannot be replaced by the batch dimension,
because positions inside a sequence have different prefixes and are different supervised
decisions.

Teacher forcing also avoids running every prefix separately. For a sequence of length
$T=4096$, separate prefix forwards would process
\begin{equation}
  1+2+\cdots+4096=8{,}390{,}656
\end{equation}
token positions through token-wise projections and MLPs. One teacher-forced forward pass
processes 4096 token positions, while the causal attention operation produces all allowed
prefix interactions together.

A training iteration is therefore:
\begin{enumerate}
  \item sample and tokenize a batch, producing IDs of shape $[B,T]$;
  \item run one causal forward pass, producing logits of shape $[B,T,V]$;
  \item shift the targets by one position and compute cross-entropy;
  \item backpropagate through every block;
  \item use an optimizer to update the parameters.
\end{enumerate}

If the model has $P$ parameters and processes $N=BT$ tokens, a useful dense-model rule of
thumb is about $2PN$ FLOPs for the forward pass and about $6PN$ FLOPs for forward plus
backward. It is an approximation: the explicit $T^2$ attention term, vocabulary head,
activation recomputation, sparsity, and hardware utilization can all matter.

Training memory has several owners. Parameters require $Ps_w$ bytes and gradients require
approximately $Ps_g$ bytes. Adam-like optimization keeps two moment estimates, commonly
in 32-bit precision, adding about $8P$ bytes; mixed-precision training may also retain a
32-bit master copy of the weights. Activations from the forward pass must remain available
for backpropagation and grow roughly with $LBTD$, plus any attention matrices that the
kernel materializes. Activation checkpointing saves memory by discarding selected
activations and recomputing them during backward.
